\documentclass[11pt,a4paper]{article}
\usepackage[margin=2.3cm]{geometry}
\usepackage{amsmath,amssymb,bm}
\usepackage{booktabs}
\usepackage{hyperref}

\newcommand{\ve}[1]{\bm{#1}}
\newcommand{\er}{\ve e_\rho}
\newcommand{\ep}{\ve e_\phi}
\newcommand{\ez}{\ve e_z}
\newcommand{\ex}{\ve e_x}

\title{Fresnel and polarization compensation in a stigmatic diopter\\[2pt]
\large Model behind notebook 04}
\date{29 September 2026}

\begin{document}
\maketitle

\begin{abstract}
A Cartesian oval $\Sigma$ images an axial point without aberration, yet local
Fresnel transmission, the tilt of the $p$ field and the ray-tube amplitude
variation still distort the vector focus. We derive how a transmitted field on
$\Sigma$ maps to the angular spectrum of the image space, why that map is the
stationary-phase (Debye) limit of an exact surface-to-spectrum integral, and how
dividing the entrance field by the resulting $p$ and $s$ transfer weights makes
the transverse focal field exactly the Airy pattern. No aplanatic objective,
sine condition or Richards--Wolf apodization is assumed.
\end{abstract}

\section{Geometry}

The object point $P_0=(0,0,z_0)$ lies in a medium of index $n_o$, and its image
$F=(0,0,z_i)$ in a medium of index $n_i$. The surface $\Sigma$ is the
vertex-connected Cartesian oval (GOTS form) with the vertex at the origin:
\begin{equation}
 n_o\,|Q-P_0| + n_i\,|F-Q| = n_o|z_0| + n_i z_i \qquad \forall\,Q\in\Sigma .
 \label{eq:oval}
\end{equation}
Every ray from $P_0$ therefore refracts towards $F$, and every ray reaches $F$ with
the same optical path. Notebook~04 uses

\begin{center}
\begin{tabular}{ll}
\toprule
Quantity & Value \\
\midrule
Object conjugate $z_0$ & $-10$ mm (air, $n_o=1$) \\
Image conjugate $z_i$ & $6$ mm (fused silica, $n_i=1.46$, inferred) \\
Wavelength $\lambda_0$ & 532 nm \\
Image-space NA $=n_i\sin\theta_{\max}$ & 0.84096 \\
Clear semi-aperture $a$ (surface radius) & 2.29106 mm \\
Sag at the rim & 2.7486 mm \\
Largest incidence angle on $\Sigma$ & $88.57^\circ$ \\
\bottomrule
\end{tabular}
\end{center}

Points of $\Sigma$ are labelled by their cylindrical radius $r$ and azimuth
$\phi$; $\er,\ep,\ez$ is the local cylindrical basis. For each $Q(r,\phi)$ let
$\ve n$ be the unit normal, $\ve d_o=(Q-P_0)/|Q-P_0|$ the incident direction,
$\ve d_i$ the Snell transmitted direction, $L_i=|F-Q|$ and
$\cos\theta=\ve d_i\cdot\ez$. Stigmatism means $\ve d_i=(F-Q)/L_i$; the notebook
verifies this to $6\times10^{-16}$ and the constancy of~\eqref{eq:oval} to
$2\times10^{-11}\lambda_0$.

\section{Entrance field and local transmission}

\paragraph{Entrance field.} With $\tan a_0=r/(z_Q-z_0)$ the tilt of the incident
ray, an x-polarized field made transverse to $\ve d_o$, with independent $p$ and
$s$ amplitudes $a_p(r)$, $a_s(r)$, is
\begin{equation}
 \ve E_0 = a_p\cos\phi\,(\er-\tan a_0\,\ez) - a_s\sin\phi\,\ep .
 \label{eq:entrance}
\end{equation}
Indeed $\ve d_o=\sin a_0\,\er+\cos a_0\,\ez$ gives $\ve E_0\cdot\ve d_o=0$. For
$a_p=a_s=1$ the transverse components are constant, $E_{0x}=1$, $E_{0y}=0$: this
is the \emph{uncompensated} illumination.

\paragraph{Local Fresnel transmission.} With $\ve s=\ve d_o\times\ve n/|\cdot|$,
$\ve p_o=\ve s\times\ve d_o$ and $\ve p_i=\ve s\times\ve d_i$,
\begin{equation}
 \ve E_t(Q) = \mathsf T\,\ve E_0,\qquad
 \mathsf T = t_p\,\ve p_i\otimes\ve p_o + t_s\,\ve s\otimes\ve s ,
 \label{eq:T}
\end{equation}
with the electric Fresnel coefficients
$t_s=2n_o c_o/(n_o c_o+n_i c_i)$, $t_p=2n_o c_o/(n_i c_o+n_o c_i)$,
$c_o=\ve d_o\cdot\ve n$, $c_i=\ve d_i\cdot\ve n$. This is the library operator
\texttt{fresnel\_transmission\_matrix}; it is applied pointwise with the true
normal of the oval. No propagation has happened yet: $\ve E_t(Q)$ is the
transmitted field on $\Sigma$.

\section{From the field on $\Sigma$ to the focal angular spectrum}

We write the image-space field as a superposition of plane waves,
\begin{equation}
 \ve E(\ve r)=\int_{|\boldsymbol\kappa|\le k_0\mathrm{NA}}
 \ve A(\boldsymbol\kappa)\,
 e^{i[\boldsymbol\kappa\cdot\boldsymbol\rho+k_z(z-z_i)]}\,d^2\kappa,
 \qquad k_z=\sqrt{k_i^2-\kappa^2},\ k_i=n_ik_0 ,
 \label{eq:asm}
\end{equation}
with $\boldsymbol\rho$ measured from the axis. The task is to obtain
$\ve A(\boldsymbol\kappa)$ from $\ve E_t$ on $\Sigma$.

\subsection{Exact statement}

The field just behind $\Sigma$ has the local form
$\ve E(\ve x)\simeq\ve E_t(Q)\,e^{ik_0\Phi(\ve x)}$, where $\Phi$ is the optical
path from $P_0$. Its angular spectrum is an integral over $\Sigma$ (for example
the Stratton--Chu or Kirchhoff representation followed by a Fourier transform,
or equivalently our operator $Q^{-1}$ followed by $R$):
\begin{equation}
 \ve A(\boldsymbol\kappa)=\int_\Sigma \mathsf K(\boldsymbol\kappa,Q)\,
 \ve E_t(Q)\,
 e^{\,i[k_0\Phi(Q)-\boldsymbol\kappa\cdot\boldsymbol\rho_Q-k_z(z_Q-z_i)]}\,dA_Q ,
 \label{eq:exact}
\end{equation}
with a slowly varying kernel $\mathsf K$. Over the $4.6$~mm aperture the phase
spans thousands of wavelengths, so~\eqref{eq:exact} is dominated by its
stationary points.

\subsection{Stationary phase}

The phase in~\eqref{eq:exact} is stationary on $\Sigma$ where the transmitted
wavevector $k_i\ve d_i$ has transverse part $\boldsymbol\kappa$. Because the
diopter is stigmatic, the transmitted wavefront is a sphere centred on $F$ and
there is exactly one such point $Q(\boldsymbol\kappa)$ for each $\boldsymbol\kappa$
in the NA disk. Along each ray of a converging spherical wave the amplitude falls
as $1/(\text{distance to }F)$, so the field near $\Sigma$ is
\begin{equation}
 \ve E\simeq \frac{\ve a(\ve d_i)}{L}\,e^{-ik_iL}\,e^{ik_0\Phi_F},
 \qquad \ve a(\ve d_i)=L_i\,\ve E_t\bigl(Q(\ve d_i)\bigr),
 \label{eq:sphere}
\end{equation}
with $L$ the distance to $F$ and $\Phi_F$ the constant optical path to the
focus. This is the Debye setting: a spherical wave of angular amplitude $\ve a$
converging on $F$. Its field near $F$ is
\begin{equation}
 \ve E(\ve r)=-\frac{ik_i}{2\pi}\int\ve a(\ve s)\,e^{ik_i\ve s\cdot(\ve r-F)}\,d\Omega ,
 \qquad d\Omega=\frac{d^2\kappa}{k_i^2\cos\theta},
\end{equation}
and comparison with~\eqref{eq:asm} gives, up to the constant phase $-i$ common to
all components,
\begin{equation}
 \boxed{\;\ve A(\boldsymbol\kappa)=
 \frac{L_i}{2\pi\,k_i\cos\theta}\;\ve E_t\bigl(Q(\boldsymbol\kappa)\bigr)\;}
 \label{eq:map}
\end{equation}
The Hessian of the phase in~\eqref{eq:exact} at the stationary point produces the
same factor $L_i/\cos\theta$; the two derivations agree.

\subsection{Energy conservation}

Map~\eqref{eq:map} conserves power patch by patch. The power leaving
$dA$ is $\tfrac{n_i}{2Z_0}|\ve E_t|^2(\ve d_i\cdot\ve n)\,dA$. That patch
subtends $d\Omega=(\ve d_i\cdot\ve n)\,dA/L_i^2$ at $F$, and
$d^2\kappa=k_i^2\cos\theta\,d\Omega$. The spectral power of the same patch is
\[
 \tfrac{n_i}{2Z_0}(2\pi)^2|\ve A|^2\cos\theta\,d^2\kappa
 =\tfrac{n_i}{2Z_0}(2\pi)^2\frac{L_i^2|\ve E_t|^2}{(2\pi k_i\cos\theta)^2}
   \cos\theta\,k_i^2\cos\theta\,d\Omega
 =\tfrac{n_i}{2Z_0}|\ve E_t|^2(\ve d_i\cdot\ve n)\,dA .
\]
The notebook checks the integrated balance to better than $10^{-9}$.

\subsection{What is neglected}

Stationary phase keeps the ray contribution of each surface point and drops
\begin{itemize}
 \item diffraction from the aperture rim, of relative size
 $\sim N_F^{-1/2}$ with Fresnel number
 $N_F=n_i a^2/(\lambda_0 z_i)\approx 2.4\times10^3$;
 \item curvature corrections of order $1/(k_iL_i)\sim10^{-5}$;
 \item multiple reflections and any field not launched from $P_0$.
\end{itemize}
It is therefore the asymptotic limit of $Q^{-1}$ followed by the exact angular
spectrum, not the discrete operator itself. The lattice operators
$Q\to T\to Q^{-1}\to R$ would require about $10^4$ samples per axis across the
aperture and are not used here. A direct test would compare both routes on a
scaled-down diopter while $N_F$ grows.

\section{Azimuthal structure and transfer weights}

For the entrance field~\eqref{eq:entrance}, rotational symmetry of $\Sigma$ fixes
the dependence on the spectral azimuth $\psi$
($\boldsymbol\kappa=\kappa(\cos\psi,\sin\psi)$; the corresponding surface point
has $\phi=\psi+\pi$):
\begin{equation}
 A_x=c_0+c_2\cos2\psi,\qquad A_y=c_2\sin2\psi,\qquad A_z=c_z\cos\psi ,
 \label{eq:azimuth}
\end{equation}
with $c_0,c_2,c_z$ functions of $\kappa$. The notebook confirms this form at an
independent azimuth to $10^{-12}$.

The two channels are characterised by the transfer weights
\begin{equation}
 w_p(r)=A_x\big|_{a_p=1,\,a_s=0,\,\psi=0},\qquad
 w_s(r)=A_x\big|_{a_p=0,\,a_s=1,\,\psi=\pi/2}.
\end{equation}
Through~\eqref{eq:entrance}, \eqref{eq:T} and~\eqref{eq:map} they contain the
Fresnel coefficients $t_p,t_s$, the entrance obliquity, the tilt of $\ve p_i$ out
of the transverse plane, and the ray-tube factor $L_i/\cos\theta$. For general
$a_p,a_s$ the transverse spectrum is
\begin{equation}
 \ve A_\perp=a_p w_p\cos\psi\,\hat{\boldsymbol\kappa}-a_s w_s\sin\psi\,\hat{\boldsymbol\psi},
\end{equation}
so $c_0=\tfrac12(a_pw_p+a_sw_s)$ and $c_2=\tfrac12(a_pw_p-a_sw_s)$. Both weights
fall towards the rim, and $w_s$ falls faster; with $a_p=a_s=1$ this makes
$c_2\neq0$ and $c_0$ non-uniform.

\section{Compensation}

Choose
\begin{equation}
 a_p=C\,\frac{w_p(0)}{w_p(r)},\qquad a_s=C\,\frac{w_s(0)}{w_s(r)} .
\end{equation}
Then $a_pw_p=a_sw_s=C\,w_p(0)$ for every $r$ ($w_p(0)=w_s(0)$ on axis), so
$c_2=0$ and $\ve A_\perp=C\,w_p(0)\,\ex$ is uniform over the NA disk. The constant
$C$ equalises the incident power through $\Sigma$,
$P_{\rm in}=\tfrac{n_o}{2Z_0}\int_\Sigma|\ve E_0|^2(\ve d_o\cdot\ve n)\,dA$, with that
of the uncompensated field. The longitudinal part $A_z$ follows from
transversality and is not prescribed.

\section{Focal field}

In the focal plane $z=z_i$ the $\psi$ integrals of~\eqref{eq:asm} with
\eqref{eq:azimuth} follow from
$\int_0^{2\pi}\cos(m\psi)\,e^{ix\cos(\psi-\varphi)}d\psi=2\pi i^mJ_m(x)\cos m\varphi$:
\begin{align}
 E_x&=F_0-F_2\cos2\varphi, & E_y&=-F_2\sin2\varphi, & E_z&=iF_z\cos\varphi,\\
 F_0&=2\pi\!\int_0^{k_0\mathrm{NA}}\!c_0J_0(\kappa\rho)\,\kappa\,d\kappa, &
 F_2&=2\pi\!\int_0^{k_0\mathrm{NA}}\!c_2J_2(\kappa\rho)\,\kappa\,d\kappa, &
 F_z&=2\pi\!\int_0^{k_0\mathrm{NA}}\!c_zJ_1(\kappa\rho)\,\kappa\,d\kappa .
\end{align}
The integrals are evaluated with Gauss--Legendre quadrature (400 nodes; doubling
changes them by less than $10^{-10}$ relative).

For the compensated field $c_2=0$ and $c_0$ is constant, so
\begin{equation}
 \ve E_\perp=2\pi c_0\,(k_0\mathrm{NA})^2\,\frac{J_1(v)}{v}\,\ex,
 \qquad v=k_0\,\mathrm{NA}\,\rho ,
\end{equation}
the Airy amplitude, with first zero at
$r_A=3.8317/(k_0\mathrm{NA})=0.61\,\lambda_0/\mathrm{NA}$.

\section{Results of notebook 04}

\begin{center}
\begin{tabular}{lcc}
\toprule
Quantity & Uncompensated & Compensated \\
\midrule
First transverse minimum along $x$ ($\mu$m) & 0.5177 & 0.3858 \\
First transverse minimum along $y$ ($\mu$m) & 0.5471 & 0.3858 \\
Airy first-zero radius $r_A$ ($\mu$m) & 0.38579 & 0.38579 \\
Departure of $\ve E_\perp$ from the Airy focus & 47.2\,\% & $<10^{-4}$\,\% \\
$\max|E_x/E_x(0)-2J_1(v)/v|$ & -- & $4\times10^{-14}$ \\
Transmitted / incident power at $\Sigma$ & 0.858 & 0.537 \\
\bottomrule
\end{tabular}
\end{center}

The compensation needs entrance amplitudes of 8--9 times the axial value near the
rim, where incidence on $\Sigma$ approaches grazing ($88.6^\circ$) and both
transfer weights fall to about 5\,\% of their axial value. This is why the
recovered Airy focus costs throughput at equal incident power.

\end{document}
